% ============================ 三、模型假设与符号说明 ============================
\section{模型假设与符号说明}

\subsection{模型假设}

\begin{enumerate}[label=\textbf{假设\arabic*.},leftmargin=3.4em,itemsep=2pt]
\item \textbf{样品代表性假设}：250 款商业啤酒构成比利时商业啤酒市场的代表性样本，
      其风格分布与化合物分布可代表该类产品总体。该假设仅用于把样本结论外推到
      “同类商业啤酒”，不外推到工业拉格或家庭自酿。
\item \textbf{品评一致性假设}：同一款啤酒的 0--6 分感官分是训练有素品评小组的
      一致估计，评分的随机误差零均值且与化合物浓度独立。该假设使回归的目标函数
      无系统偏差，仅含方差。
\item \textbf{对数变换假设}：\texttt{1.csv} 中未标注量纲的多数列是浓度的
      $\log_{10}$ 变换（负值对应浓度 $<1$），少数列为原始浓度；具体到哪一列由
      \S\ref{sec:logbase} 用校准浓度表逐列反证，不作整体假定。
\item \textbf{检出限假设}：某化合物在大量样本上取到同一最小值时，该值是仪器检出
      限而非真实浓度，真实值落在区间 $(-\infty,\ v_{\min}]$；建模时以“封底
      指示变量”显式吸收这一信息。
\item \textbf{风格可分假设}：\texttt{tasting\_category\_fine} 给出的风格标签在
      化学谱上可分，即存在从化合物谱到风格的判别函数，且其准确率显著高于随机
      基线。该假设是问题三“风格条件生成”可评估的前提，并由
      \S\ref{sec:q3style} 实测检验。
\item \textbf{迁移可行性假设}：跨风格迁移的目标化合物谱落在训练样本张成的化学
      可行域内，因此沿潜空间直线插值不会产生化学上不可能的负浓度或超量级浓度；
      该假设由迁移后浓度的边界检查实测核验。
\end{enumerate}

\subsection{符号说明}

本文全局符号列于表~\ref{tab:symbols}，各章局部符号在首次出现处另行说明。

\begin{table}[H]
\centering
\caption{主要符号说明}
\label{tab:symbols}
\small
\begin{fittable}
\begin{tabular}{@{}lll@{}}
\toprule
符号 & 含义 & 单位 / 取值域 \\
\midrule
$n$ & 啤酒样本数 & 250 \\
$p$ & 化合物特征维数 & 231 \\
$q$ & 感官维度数 & 50 \\
$\bm{X}\in\R^{n\times p}$ & 化合物丰度矩阵（插补后） & 无量纲 \\
$x_{ij}$ & 第 $i$ 款酒第 $j$ 个化合物的数值 & 见 \S\ref{sec:logbase} \\
$y_{ik}$ & 第 $i$ 款酒第 $k$ 个感官维度得分 & 0--6 分制 \\
$s_i$ & 第 $i$ 款酒的风格标签 & 22 类之一 \\
$\bm{z}_i\in\R^{K}$ & 第 $i$ 款酒在 $K$ 维潜因子空间的坐标 & 无量纲 \\
$K$ & 潜因子个数（并行分析确定） & 正整数 \\
$\rho_{jk}$ & 化合物 $j$ 与 $k$ 的皮尔逊相关系数 & $[-1,1]$ \\
$w_{jk}$ & 共表达图边权，$w_{jk}=|\rho_{jk}|$ & $[0,1]$ \\
$\bm{A},\bm{D}$ & 图邻接矩阵与度矩阵 & --- \\
$\bm{S}$ & 图扩散算子，$\bm{S}=(1-\alpha)\bm{I}+\alpha\tilde{\bm{A}}$ & --- \\
$R^2$ & 决定系数，$1-\sum(y-\hat y)^2/\sum(y-\bar y)^2$ & $(-\infty,1]$ \\
$\Phi^{-1}$ & 标准正态分位数函数 & --- \\
$T_c$ & 化合物 $c$ 的气味阈值 & $\mu\mathrm{g\cdot L^{-1}}$ 或 $\mathrm{mg\cdot L^{-1}}$ \\
$\mathrm{OAV}_c$ & 气味活度值，$\mathrm{OAV}_c=C_c/T_c$ & 无量纲 \\
$\mathrm{FD}$ & Fréchet 距离（分布层面的生成质量） & 同 $\bm{X}$ 量纲 \\
$\lambda$ & 跨风格迁移强度，$0\le\lambda\le1$ & 无量纲 \\
$B$ & bootstrap 重采样次数 & 200 \\
$\alpha$ & 图扩散强度（正则参数） & $[0,1]$ \\
\bottomrule
\end{tabular}
\end{fittable}
\end{table}
